\section{Triple ramification and multisections of a spin line}\label{p8-sec:graphs}

The triple-ramification divisors give a test for the appearance of
actual curves. A multisection here is the image of a curve in the
total space of a line or vector bundle whose projection to the base
is finite and surjective. We call its parametrization simple if it
has generic degree one onto its image.

\begin{theorem}\label{p8-thm:algebraic-graphs}
Let $C$ be any smooth projective curve of genus $g\ge2$, and let
$\theta^2=K_C$. On every connected component of the space of connected
degree-$m$ covers of $C$ with exactly one triple ramification point
and no other ramification, there is a cover $f:Y\to C$ such that
\[
                       h^0(Y,f^*\theta)>h^0(C,\theta).
\]
For every $m\ge3$ the holomorphic total space of $\theta$ consequently
contains simple multisections of degree $m$ with smooth normalization
of genus $H_m+1$ and a single triple ramification point over $C$.
Their supports can be chosen arbitrarily close to the zero section.
\end{theorem}
\begin{proof}
Let $T$ be such a connected component. It is a proper curve finite
in coarse space over $C$ by Section~\ref{p8-sec:hurwitz}. We may pass
to a connected finite level cover so that the family
$\pi:\mathcal Y\to T$ is a smooth projective family over a smooth
curve. For example, a symplectic level structure of level at least
three on the first homology of the smooth sources rigidifies their
automorphisms. Write $f:\mathcal Y\to C$ also for the evaluation map,
$p:T\to\mathcal Y$ for the ramification section, $E=p(T)$, and
$K=c_1(f^*K_C)$. The number $K\cdot E$ is strictly positive since
$T\to C$ moves the branch value. Relative Riemann--Hurwitz and the
first nonzero jet along $p$ give
\[
 c_1(\omega_\pi)=K+2E,\qquad K^2=0,\qquad E^2=-K\cdot E/3.
\]
Put $L=f^*\theta$. Grothendieck--Riemann--Roch for this smooth family
gives the degree of its negative determinant of cohomology as
\begin{align*}
 \deg c_1(-R\pi_*L)
 &=-\int_{\mathcal Y}
    \left(\tfrac12c_1(L)^2
       -\tfrac12 c_1(L)c_1(\omega_\pi)
       +\tfrac1{12}c_1(\omega_\pi)^2\right)\\
 &=\tfrac16 K\cdot E-\tfrac13 E^2=\tfrac5{18}K\cdot E .
\end{align*}

Suppose that every section on every fiber is a pullback. By cohomology
and base change $\pi_*L$ is then the constant bundle with fiber
$H^0(C,\theta)$. Relative duality and Riemann--Roch show that
$S=\pi_*(\omega_\pi L^{-1})/(df)\pi_* L$ is a line bundle, with
$\deg S^*=\deg c_1(-R\pi_*L)$. By the residue pairing $S$ is
isotropic in the rank-two orthogonal residue bundle. That bundle has
the global isotropic line $I=\pi_*(L(E)/L)$ of degree $K\cdot E/6$.
The two isotropic directions in an orthogonal plane are distinct; once
one is chosen, the other is a line paired perfectly with it.
Connectedness therefore forces $\deg S^*=\pm K\cdot E/6$.
Both values contradict the degree $5K\cdot E/18$. This argument
uses only the algebraic isotropics and includes effective $\theta$.

For each $m\ge3$ choose standard surface generators. On the first
handle use $(12),(23)$, and on the second use $(1\,2\,\ldots\,m)$
and the identity, with the remaining generators trivial. The
commutator product is a $3$-cycle, and these permutations generate
$S_m$. Filling the puncture gives a connected cover with the required
ramification. Its Hurwitz component consists of covers with the same
full monodromy. Take the additional section just proved to exist.
Its graph $(f,s)$ is simple: the normalization of its image gives an
intermediate cover, whereas the natural action of $S_m$ on $m$
sheets is primitive. A normalized image of degree one over $C$
would make $s$ a pullback, and the other possibility makes the graph
generically injective. The graphs of positive scalar multiples of
$s$ have different maximum fiber norms and converge to the zero
section. Their normalizations have the stated genus by
Riemann--Hurwitz.
\end{proof}

\begin{theorem}\label{p8-thm:path-graphs}
Let $a,b$ be distinct non-effective theta characteristics on $C$.
Identify their underlying smooth complex lines, and take any
continuous path of real Cauchy--Riemann operators on this line
between their spin operators. Suppose the operators on $C$ are
invertible throughout. In every degree $m\ge3$ there is an interior
parameter at which a smooth cover $f:Y\to C$ with full symmetric
monodromy and a single triple ramification point has nonzero
pullback kernel. Its kernel sections give simple almost-complex
multisections accumulating at the zero section. Such paths and
pairs exist in every base genus $g\ge2$. A single split path on
a closed symplectic six-torus realizes these multisections in all
degrees $m\ge3$.
\end{theorem}
\begin{proof}
Choose the first generator in the preceding permutation construction
to represent the Poincar\'e dual of the nonzero difference $ba^{-1}$
modulo two. Its monodromy is odd. On the corresponding compact
triple divisor $T$, the ratio of the two induced spin parities is
$-1$ by Section~\ref{p8-sec:characters}, since both embedded spin lines
have even parity. By Theorem~\ref{p8-thm:profiles} the oriented
obstruction planes at the two endpoints have opposite nonzero Euler
degrees on $T$. If every pullback along the path were injective
there, the constant-rank smooth-domain Fredholm family would give
an oriented plane bundle over $T\times[0,1]$. Its two Euler degrees
would be equal. This proves the existence of a nonzero section.
It cannot descend since the base operator has zero kernel. In using
primitivity for its graph, normalization of a simple image is the
usual factorization for a nonconstant pseudoholomorphic map.
The projection from the normalization to $C$ is holomorphic.
A degree-one projection is an isomorphism, and the resulting graph
over $C$ solves the base equation. This proves simplicity just as
in the holomorphic case. The positive real multiples again give
distinct supports. The assertion permits the same parameter to
serve more than one degree.

To construct a path, decrease the zeroth-order term of the spin
operator on $a$ to zero. Every strictly positive amplitude is
injective by the spin norm identity, and $\bar\partial_a$ is
invertible by non-effectivity and index zero. Join the complex
operators by a path in the complement of the theta divisor of
$\operatorname{Pic}^{g-1} C$, which is connected. A family of
representatives along the interval identifies the underlying
smooth lines. Reverse the amplitude deformation at $b$.
For existence in each genus take a hyperelliptic curve
$y^2=\prod_{i=1}^{2g+2}(x-e_i)$, with ramification points $W_i$
and divisor $H$ over infinity. For $|A|=g+1$ the line
$\mathcal O_C(\sum_{i\in A}W_i-H)$ squares to $K_C$ and has no
sections. Indeed, an invariant meromorphic function with at most
these simple poles and vanishing over infinity is zero. An
anti-invariant one is $yP(x)/\prod_{i\in A}(x-e_i)$ for a polynomial
$P$, and vanishing over infinity again forces zero. Replacing one
element of $A$ changes the line, since two distinct points on $C$
are not linearly equivalent. Finally keep one spin summand fixed
and put the path in the other summand of the normal bundle of
degree $2g-2$. The graph equation defines a path of almost-complex
germs with the prescribed projections. The compact realization
of \cite[Proposition~3.1 and \S15]{WueUnramified} extends the family from a fixed small
tube to tame structures on a symplectic six-torus. For each kernel
section a sufficiently small multiple lies in this tube, giving
a curve in the class $m[C]$ with the specified normalization.
\end{proof}
